\section{Búsqueda + LLM: la aplicación estrella de las empresas}

\begin{importantbox}{Grounding with Search}
Casi todos los casos de uso empresariales combinan un LLM con la búsqueda:
\begin{itemize}
    \item El conocimiento del LLM tiene un límite de vigencia temporal: se necesita la búsqueda para obtener información actualizada
    \item El LLM produce alucinaciones: se necesita la búsqueda para verificar los hechos
    \item El LLM es bueno para razonar pero carece de referencias: la búsqueda aporta citas precisas de las fuentes
\end{itemize}
Esto hace que RAG (generación aumentada por recuperación) se haya vuelto la arquitectura estándar de las aplicaciones empresariales de IA generativa.
\end{importantbox}

Desde el punto de vista de la ingeniería, el pipeline de RAG debe resolver tres problemas clave: primero, la latencia que trae consigo el crecimiento del índice vectorial; segundo, las anomalías de recuperación (que se devuelvan documentos obsoletos o erróneos); tercero, la estrategia de integración entre el LLM y los resultados recuperados. Burak comentó que algunos equipos agregan una capa de «fact filter» antes de los resultados de la recuperación: primero se le entregan los resultados a un smaller model para hacer un reordenamiento fino, y después se le entregan al LLM grande para la generación.

\subsection{Resumen de la sección}

Búsqueda + LLM es el patrón de aplicación de IA generativa más común en las empresas, y resuelve los tres puntos críticos de la vigencia temporal, las alucinaciones y las referencias.

